\ssr{Введение}

Цели данной лабораторной работы:
\begin{enumerate}
	\item изучение системных функций выделения памяти в ОС GNU/LINUX;
	\item получение практических навыков работы с динамической памятьюв ОС GNU/LINUX.	
\end{enumerate}

\ssr{Выполнение}

\ssection{Задание 1. Программа}

Программа на C/C++ работет в одном из 2-х режимов, режим работы задается при запуске в качестве аргумента командной строки:
\begin{itemize}
	\item Режим 1 (Р1). Сбалансированное выделение/освобождение памяти;
	\item Режим 2 (Р2). Контролируемая "утечка" памяти.
\end{itemize}

\lstinputlisting[
	style=cstyle, 
	caption={Программа с двумя режимами}
]{program.c}

\begin{figure}[!htbp]
	\centering
	%\includegraphics[width=0.6\linewidth]{""}
	\caption{Демонстрация работы}
\end{figure}
\FloatBarrier

\ssection{Задание 2. bash-скрипт}

\lstinputlisting[
	style=bash, 
	caption={bash-скрипт для автоматизации сбора информации о виртуальной памяти целевого процесса}
]{vmem.sh}

\ssection{Задание 3. Автозапуск скрипта по таймеру}

\begin{figure}[!htbp]
	\centering
	%\includegraphics[width=0.6\linewidth]{""}
	\caption{Файлы с картами памяти}
\end{figure}
\FloatBarrier

\ssection{Дополнительное задание. Анализ карт памяти}

\begin{figure}[!htbp]
	\centering
	%\includegraphics[width=0.6\linewidth]{""}
	\caption{График изменения кучи и RSS по времени}
\end{figure}
\FloatBarrier

\ssr{Вывод}

В ходе лабораторной работы были:  
\begin{itemize}
	\item изучены системные функции выделения памяти в ОС GNU/LINUX;
	\item получены практические навыки работы с динамической памятьюв ОС GNU/LINUX.	
\end{itemize}

\ssr{Контрольные вопросы}

\begin{enumerate}
	\item \textit{Что такое страница памяти и какие размеры страниц поддерживаются в Linux?}\newline 
	Блок фиксированного размера, являющийся единицей управления памятью (выделения, отображения файлов). 
	В Linux типичный размер составляет 4 КБ, но архитектура может поддерживать и другие размеры (можно посмотреть через getpagesize()).
	\item \textit{Какие существуют типы выделения памяти в программах на C?}\newline 
	Статическое выделение (глобальные/статические переменные, определяемые при компиляции), автоматическое выделение (локальные переменные в стеке), динамическое выделение (куча, через malloc/calloc/realloc).
	\item \textit{Что такое виртуальная память процесса и какие основные сегменты она включает (stack, heap, data, text)?}\newline 
	Абстрактное адресное пространство, предоставляемое ядром каждому процессу, основная структура которого включает: 
	text (машинный код), 
	data (инициализированные глобальные/статические переменные), 
	bss (неинициализированные глобальные/статические переменные), 
	heap (динамическое выделение, растёт в сторону высоких адресов), 
	stack (локальные переменные и вызовы функций, растет в сторону низких адресов).
	\item \textit{Как функции malloc() и free() взаимодействуют с ядром ОС?}\newline 
	malloc при нехватке памяти в куче запрашивает у ядра увеличение границы кучи через brk/sbrk (для небольших выделений) или отображает новую область через mmap (для больших выделений).
	free обычно лишь возвращает блок в свободный список внутри библиотеки C и не возвращает память ядру немедленно, поэтому системный вызов происходит только тогда, когда библиотека решает сократить кучу .
	\item \textit{Всегда ли вызов malloc() приводит к системному вызову?}\newline 
	Не каждый вызов malloc инициирует системный вызов: библиотека C заранее получает у ядра крупные блоки памяти, а последующие небольшие выделения выполняются из уже имеющегося пула, что снижает накладные расходы на переход в режим ядра.
	\item \textit{Что показывает карта виртуальной памяти процесса и как идентифицировать в ней кучу (heap)?}\newline 
	Карта виртуальной памяти (/proc/pid/maps) показывает области отображения адресного пространства процесса (начальный/конечный адрес, права доступа, источник отображения). 
	Кучу можно идентифицировать по: анонимному отображению с правами rw-p и меткой [heap] (в старых ядрах/раскладке через brk) либо по большому количеству анонимных rw-p областей, созданных через mmap (в современных glibc при многопоточности/крупных выделениях).
	\item \textit{Почему размер кучи в картах памяти может увеличиваться, но не уменьшаться даже после вызовов free()?}\newline 
	Размер кучи не уменьшается: free обычно лишь помечает блок как свободный и возвращает его во внутренний свободный список для повторного использования, но не возвращает страницы ядру.
	\item \textit{Какие метрики памяти наиболее важны для анализа при исследовании работы программы с динамической памятью?}\newline 
	RSS (резидентный набор, фактически занятые физические страницы), 
	VSZ (размер виртуального адресного пространства), 
	размер кучи (выделенная программой динамическая память), количество страничных ошибок, 
	PSS (пропорциональный резидентный набор, учитывающий общие страницы). 
	При анализе утечек особое внимание следует уделять тренду роста RSS и размера кучи.
	\item \textit{Как отличить штатный рост потребления памяти от утечки памяти при анализе серии карт памяти?}\newline 
	При анализе нескольких снимков карты памяти штатный рост проявляется как ступенчатое увеличение с последующей стабилизацией (например, прогрев кэша, предварительное выделение пула).
	Утечка проявляется как непрерывный монотонный рост без стабилизации, и после многократного выполнения одного и того же сценария (например, открытие/закрытие, запрос/освобождение) занятая память не возвращается к исходному уровню .
	\item \textit{Почему Resident Set Size (RSS) может отличаться от размера кучи в картах памяти?}\newline 
	Куча -- это логическая занятая память, заявленная malloc в адресном пространстве, но не все страницы кучи физически загружены (например, выделенные, но не записанные страницы не занимают физическую память). 
	Кроме того, RSS включает стек, сегменты кода, разделяемые библиотеки, отображённые файлы и другие части, поэтому RSS обычно больше фактически используемой кучи, но может быть и меньше виртуально выделенной кучи.
\end{enumerate}
